% commentaire.tex — f176r, vv. 1128–1168 (conversion initiale des travaux antérieurs)
\comm{1127}{τὸ ξιπομάκαιραν.}{La harpé de Persée. (Note de Willems.)}
\comm{1130}{σκαιοῖσι γάρ τοι καινὰ προσφέρων σοφὰ.}{Cf. le reproche qu'Euripide mit dans la bouche de Médée dans la tragédie de ce nom (vers 298) :\citComm{σκαιοῖσι μὲν γὰρ καινὰ προσφέρων σοφὰ.}}
\comm{1133}{οἷον ἐπιτηκίζει μοι}{$=$ οἷον ἐπιθήκιζέ μοι. À propos des images animales du trompeur : du renard, ἡ ἀλώπηξ, et du singe, ὁ πίθηκος, vus comme les représentants de la ruse et de la fourbe, le grand Taillardat a consacré deux paragraphes auxquels je renvoie le lecteur : §§405-406. Notons que sur le même modèle que πιθηκίζειν (qui se rencontre aussi au vers 1290 des \italique{Guêpes}), il existe chez Aristophane le verbe ἀλωπεκίζειν (\italique{Guêpes}, vers 1241).}
\comm{1136–59}{Chant ibycéen du chœur.}{Pour la structure et le texte, voir le commentaire critique. Mis de cette façon en ordre, le chant est entrecoupé entre chaque paire strophique d'un intermède, μεσῳδός, commençant par un impératif de φαίνω : aux vers 1143-4, φάνηθ᾽ ὦ τυράννους στυγοῦσ᾽ ὥσπερ εἰκός, et au vers 1154, φαίνετον ἄμβροτον ὄψιν.}
\comm{1151}{ἵνα,}{Une glose à l'intérieur du texte du manuscrit\marqueur{} précise quel sens donner à (ἵνα : ) ὅπου\nc{On retrouve une glose identique dans de nombreux lexiques : Συναγωγή ι 82 ; Hésychius ι 648 ; Photius ι 133 ; Suidas ι 360 ; \&c.}.}
\comm{1168–9}{ἃ νῦν ὑποικουρεῖτε τοῖσιν ἀνδράσιν / ἀπὸ τῆς στρατιᾶς παροῦσιν.}{À propos de la légèreté des Athéniennes pendant le service de leur mari, cf. le vers 495 et le commentaire \italique{ad locum}.}
\comm{1229–31}{τὼ Θεσμοφόρω δ᾽ / ἡμῖν ἀγαθὴν / τούτων χάριν ἀνταποδοίτην.}{Selon son habitude -- qui était peut-être d'usage dans l'Ancienne comédie --, Aristophane conclut sa pièce par des souhaits de succès (il vise bien-sûr la victoire au concours comique).}
